\section{Embedded primary strata for connected local algebras}
\label{sec:embedded-multigenerator}
The incidence of Theorem~\ref{thm:full-codimension-two} transfers
more than the support of a quotient scheme. On its quadratic-generation
open it transfers its local rings, and hence its primary structure.
We apply it to a connected noncurvilinear family in which the full
multiplication failure scheme has an embedded component.

For $h\ge5$ set
\begin{equation}\label{eq:fork-algebra}
 B_h=\C[x,y]/(x^2,xy,y^h),\qquad \dim_\C B_h=h+1.
\end{equation}
Let $F_0=0$, $F_1=1$, and $F_{j+1}=uF_j+vF_{j-1}$, and put
\begin{equation}\label{eq:fork-ideals}
 \begin{split}
 J_h&=(F_h,vF_{h-1})\subset\C[u,v],\\
 I_h&=(F_h,vF_{h-1},bv,b^2u)\subset\C[b,u,v],\\
 P_h&=(u,v),\qquad Q_h=(b^2,bv,F_h,vF_{h-1}).
 \end{split}
\end{equation}
The equality $J_h=(F_h,F_{h+1})$ follows from the recurrence.

\begin{theorem}[A connected embedded-primary family]
\label{thm:embedded-multigenerator}
Consider the unital Grassmannian of $(h-1)$-planes in $B_h$, and
its full zeroth-Fitting multiplication failure scheme $D_h$ in
any degree $m\ge3$. At the plane
\begin{equation}\label{eq:fork-plane}
 W_0=\langle1,x,y^2,y^3,y^5,\ldots,y^{h-1}\rangle
 \quad\text{(the term $y^4$ is omitted)},
\end{equation}
there is a Zariski neighbourhood in $D_h$ isomorphic to
\begin{equation}\label{eq:fork-local-ring}
 \operatorname{Spec}\bigl((\C[b,u,v]/I_h)
                [t_1,\ldots,t_{h-2}]\bigr)_f,
 \qquad f(0)=1.
\end{equation}
Here the variables $t_i$ are actual Grassmannian coordinates, and
$f$ is the coefficient of a specified quadratic product in the
one-dimensional quotient of the conductor hyperplane by $W$.
This is a local description of the full failure scheme, not of a
transverse test family.

The ideal in \eqref{eq:fork-ideals} has the irredundant primary
 decomposition
\begin{equation}\label{eq:fork-primary}
 I_h=P_h\cap Q_h.
\end{equation}
Its associated primes are $(u,v)$ and $(b,u,v)$. The first is the
minimal prime and the second is embedded. The same decomposition,
with the ideals extended and localized, describes
\eqref{eq:fork-local-ring}. The reduced support has codimension
$h-3\ge2$ in the unital Grassmannian, so this failure scheme is not
Cartier. Its nilradical has exact nilpotency index $2h-3$.
For the displayed choice of embedded primary ideal,
\begin{equation}\label{eq:fork-lengths}
 \dim_\C\C[b,u,v]/Q_h=\binom h2+h-1,
 \qquad
 \dim_\C\sqrt{I_h}/I_h=\binom h2+h-3.
\end{equation}
The second formula concerns the coefficient ring before adjoining
the Grassmannian variables. It records a nilradical supported at
the embedded point, rather than multiplicity along the generic
point of a Cartier divisor.
\end{theorem}
\begin{proof}
We first compute an entire open chart of $\operatorname{Hilb}^2(B_h)$.
Choose the chart on which $1,y$ is a basis of the universal quotient.
Write $y^2=uy+v$ and $x=a+by$. The relation $xy=0$ gives
$a=-bu$ and $bv=0$. After this substitution, $x^2=0$ is equivalent
to $b^2u=0$: its constant equation follows from $bv=0$ and $b^2u=0$.
Finally
\[
 y^h=F_h y+vF_{h-1}
\]
shows that $y^h=0$ is equivalent to the two generators of $J_h$.
These calculations give a quotient algebra free on $1,y$ over
$\C[b,u,v]/I_h$, and conversely every such quotient has these
coefficients. They prove an isomorphism of the chart functor with
$\operatorname{Spec}\C[b,u,v]/I_h$ over arbitrary coefficient rings.
No radical has been taken.

Here is a direct primary calculation. Write
$A_h=\C[u,v]/J_h$ and $\mathfrak n=(u,v)A_h$. It is an Artin
local algebra, with residue field $\C$. In $A_h[b]$, coefficient
comparison in the variable $b$ gives
\[
 \mathfrak n A_h[b]\cap (b^2,bv)
       =(bv,b^2\mathfrak n)=(bv,b^2u).
\]
Lifting to $\C[b,u,v]$ proves \eqref{eq:fork-primary}.
The ideal $P_h$ is prime; the quotient by $Q_h$ is an Artin local
ring supported at $(b,u,v)$, so $Q_h$ is primary. The decomposition
is irredundant: $b^2\in Q_h\setminus P_h$, whereas
$u\in P_h\setminus Q_h$. The asserted associated primes follow.
Polynomial extension and localization at a function nonzero at the
specified point retain both components.

For clarity, the numerical invariants can be read without any
primary-decomposition algorithm. The ring $A_h$ is the contact
algebra $G_h$ of Theorem~\ref{thm:contact-hyperplane-primary}.
It has dimension $\binom h2$, and
\[
 A_h/(v)\simeq\C[u]/(u^{h-1}),\qquad
 \mathfrak n^{2h-3}=0,\qquad u^{2h-4}\ne0.
\]
These last assertions also follow from its two-row Schur basis:
$\deg u=1$, $\deg v=2$, the top weighted degree is $2h-4$, and the
coefficient of the top Schur class in $u^{2h-4}$ is nonzero.
This is the classical basis input recorded in \cite{Grinberg}.
As vector spaces the quotient by $Q_h$ is
$A_h\oplus b(A_h/(v))$, proving the first length formula.
In the quotient by $I_h$, the nilradical is
\[
 \mathfrak n\ \oplus\ b\bigl(u\C[u]/(u^{h-1})\bigr).
\]
All coefficients of $b^j$ for $j\ge2$ are scalars. This proves the
second length formula. The inclusion of the constant coefficient
ring $A_h$ is injective, so the upper bound $2h-3$ on nilpotency is
attained. Localization at $f(0)=1$ does not kill its nonzero top
class: it is supported at the retained maximal ideal.

It remains to prove that this chart is realized by the full
multigenerator failure scheme. At $b=u=v=0$, the conductor hyperplane
is
\[
 E_0=\langle1,x,y^2,\ldots,y^{h-1}\rangle.
\]
The plane $W_0$ in \eqref{eq:fork-plane} is a hyperplane in $E_0$,
and $W_0^2=E_0$, since $(y^2)^2=y^4$. In the universal quotient
chart the conductor hyperplane has the basis
\[
 1,\quad x-by,\quad y^j-F_jy\quad(2\le j\le h-1).
\]
Take the Grassmannian chart in which the vector $y^4-F_4y$ is the
omitted direction, and add its multiples $t_i$ to each of the other
nonconstant basis vectors. This gives $h-2$ free coordinates on
$\Gr(h-2,E/\mathcal O)$. Let $q_2$ be the resulting lift of the
$y^2$ basis vector. Its square belongs to $E$, because $E$ is an
algebra. The coefficient $f$ of $q_2^2$ in $E/W$, in the omitted
basis direction, equals $1$ at the origin. On $f\ne0$ one has
$W^2=E$. Theorem~\ref{thm:full-codimension-two} now identifies this
localized polynomial chart with an open neighbourhood of $W_0$
in $D_h$, proving \eqref{eq:fork-local-ring}.

The unital Grassmannian has dimension $2(h-2)$, whereas the support
of \eqref{eq:fork-local-ring} has dimension $1+(h-2)=h-1$.
Their difference is $h-3$. This proves the non-Cartier assertion
and finishes the proof.
\end{proof}

For $h=5$ the theorem gives four-planes in a connected length-six
algebra. Its coefficient ideal is explicitly
\[
 \bigl(u^4+3u^2v+v^2,\ u^3v+2uv^2,\ bv,\ b^2u\bigr).
\]
The reduced support is the $b$-axis, the chosen embedded primary
quotient has length $14$, and the nilradical has dimension $12$
and index $7$. Increasing $h$ changes both the higher-jet equations
and the embedded nilpotent algebra; it does not merely change the
exponent of a fixed determinant divisor. The parameter $b$ couples
the extra square-zero tangent direction $x$ to the higher contact
coefficients $u,v$ through $bv$ and $b^2u$.

\begin{corollary}[A three-generator member after adjoining a reduced factor]
\label{cor:three-generator-embedded}
For every $s\ge1$, let
$B=(\C[x,y]/(x^2,y^2))\times\C^s$. The full stable failure scheme
of unital codimension-two planes has a chart, meeting an embedded
associated stratum, of the form
\[
 \operatorname{Spec}
 \bigl((\C[a,b]/(ab,b^2))[t_1,\ldots,t_{s+1}]\bigr)_g,
 \qquad g(0)=1.
\]
Its primary decomposition is the extension of
$(ab,b^2)=(b)\cap(a,b)^2$. In particular $s=1$ gives three
multiplication generators and an embedded component on the full
failure scheme, rather than on a selected slice.
\end{corollary}
\begin{proof}
Use the open-and-closed part of $\operatorname{Hilb}^2(B)$ supported
in the first factor. Proposition~\ref{prop:noncurvilinear-embedded}
gives the coefficient chart $\C[a,b]/(ab,b^2)$ there. At its origin
the conductor hyperplane is
$\langle1,y,xy\rangle\times\C^s$. Let $e_i$ denote the standard
idempotent of the $i$th reduced factor. The plane generated by
\[
 1_B,\qquad (y,e_1),\qquad (xy,0),\qquad e_2,\ldots,e_s
\]
has rank $s+2=\dim B-2$. Its square contains $e_1=(y,e_1)^2$,
so its quadratic image is the entire conductor hyperplane.
The relative Grassmannian chart has $s+1$ parameters. Localize at
the quadratic-generation minor nonzero at this point and apply
Theorem~\ref{thm:full-codimension-two}. The displayed primary
identity is retained under this polynomial extension and localization.
\end{proof}

The connected family of Theorem~\ref{thm:embedded-multigenerator}
and the three-generator corollary serve different purposes. The
first gives unbounded connected local complexity; the second supplies
a three-generator instance of the same incidence mechanism. Both
also describe smooth local covers of the unrestricted generating
Grassmannians by Remark~\ref{rem:codim-two-unit-cover}. Neither uses
a claim that primary decomposition survives arbitrary specialization.
